\section{Am Tümpelackersee}
\stepcounter{ctrpoeminyear}
%\addcontentsline{toc}{section}{\thechapter.\arabic{ctrpoeminyear} \;}

\begin{strophe}
\vers{Es spielten zwei Jungen am Tümpelackersee.}
\vers{Sie spielten "Stöckchen, wirfs in die Höh'".}
\vers{Es dauert nicht lange bis Stock fiel ins Wasser}
\vers{und um es zu retten wird einer pitschnasser.}
\vers{Der eine, kein Schwimmer, kämpft um die Luft}
\vers{und um ihn zu retten springt auch der zweite mit Wucht.}
\vers{Jetzt beide im Wasser, sie tanzen mit Händen,}
\vers{sie stampfen, doch Wasser bändigt den ersten.}
\vers{Dieser sauft ab und fließt richtung Boden}
\vers{und um seinen besten Freund zu holen}
\vers{taucht der zweite ihm hinterher,}
\vers{doch sein Blick erhascht in diesem Meer}
\vers{eine Pforte zu einer anderen Welt.}
\vers{Er nähert sich, und ist bald umstellt}
\vers{von einer Armee an Meerjungpersonen,}
\vers{die, wie's scheint, dort unten seit jeher schon wohnen.}
\vers{Der Junge erklärt seinen Tatbestand,}
\vers{fragt nach Stöcklein, das ihm entwandt.}
\vers{Gelächter aus Reihen der Meerjungarmee}
\vers{und tanzend, wie eine himmlische Fee,}
\vers{schwimmt eine der Flossen zu ihm heran}
\vers{und lädt ihn ein auf einen Anstoß mit Wein.}
\vers{Der Junge, der nie solche Kost gewusst,}
\vers{betritt das Gebiet der Fische selbstbewusst}
\end{strophe}

\phantom{a}

\phantom{a}

\phantom{a}

\phantom{a}

\phantom{a}

\phantom{a}

\phantom{a}

\phantom{a}

\phantom{a}

\begin{strophe}
\vers{und nach fünf Jahren der Freude im Wasserland}
\vers{findet er die Idee zurückzukehren an Land.}
\vers{Dort oben ward lange gesucht nach den Zwein,}
\vers{fünfzig Jahre um genau zu sein.}
\vers{Und als der Junge seiner Mutter vors Antlitz tritt,}
\vers{schreit sie auf, als sei sie verrückt.}
\vers{Ein junger Bursche, der aussieht wie der,}
\vers{der ertrunken sein sollte im Tümpelackersee.}
\vers{In Jahrzehnten soll er nur Jährchen gealtert sein,}
\vers{und steht jetzt so vor ihr, wohl und bekleidet,}
\vers{als sei er nie weg gewesen für all die Zeiten.}
\end{strophe}

\begin{strophe}
\vers{Das ist das Ende der kurzen Geschicht,}
\vers{die auf Erden gar fünfzig Jahre hergibt.}
\vers{Die Mutter und Sohnemann lebten glücklich}
\vers{zwei Tage lang, bis die Mutter an Schwäche vor Glück}
\vers{starb und ihm alles Hab und Gut lies zurück.}
\vers{Er nahm es sodann und versank's in den Tiefen}
\vers{des Tümpelackersees, um anzubieten}
\vers{einen Dank für schöne fünf Jahre,}
\vers{beim Volk der Flossenträger hierzulande.}
\end{strophe}

13.08.2026, überarbeitet am 14.08.2026, minimal

überarbeitet am 16.08.2026

%Christoph Koloska

